% Generated by `cargo xtask refs-oracle generate` from papers.toml. Do not edit.
\documentclass[11pt]{amsart}
\usepackage{amsmath,amssymb,amsthm}
\usepackage{pdfcraftrefs}
\pagestyle{headings}
\renewcommand\labelitemi{$\bullet$}
\newtheoremstyle{pcplain}{\topsep}{\topsep}{\normalfont}{}{\itshape}{.}{.5em}{}
\theoremstyle{pcplain}
\newtheorem{theorem}{Theorem}[section]
\newtheorem{lemma}{Lemma}[section]
\newtheorem{proposition}{Proposition}[section]
\newtheorem{corollary}{Corollary}[section]
\newtheorem{claim}{Claim}[section]
\newtheorem{conjecture}{Conjecture}[section]
\newtheorem{fact}{Fact}[section]
\newtheorem*{mainthm}{Main Theorem}
\newtheorem*{utheorem}{Theorem}
\theoremstyle{definition}
\newtheorem{definition}{Definition}[section]
\newtheorem{assumption}{Assumption}[section]
\newtheorem{notation}{Notation}[section]
\newtheorem{example}{Example}[section]
\theoremstyle{remark}
\newtheorem{remark}{Remark}[section]
\newtheorem*{uremark}{Remark}
\newtheorem{observation}{Observation}[section]
\numberwithin{equation}{section}
\renewcommand\qedsymbol{$\blacksquare$}
\relpenalty=10000
\binoppenalty=10000
\sloppy
\begin{document}
\title{Spectral gaps of weighted shifts}
\author{G. Writer}
\maketitle

\pdfcraftmark{sec:a}\section{Background}\label{sec:a}


Under this assumption, the weak space $x_n$ is bounded by a symmetric operator of $u_t$. The finite solution $x_n$ is close to a smooth kernel of $\mathbb{R}^d$ whenever $f$ is controlled by $f$, and the operator is admissible. Moreover, the positive matrix $f$ is equivalent to the process $C_0$.

\pdfcraftthm{s1}\begin{notation}\label{not:a}
Every admissible kernel determines the domain $n^2$, and the constant is convex.
\end{notation}\pdfcraftthmend


\pdfcraftthm{s2}\begin{conjecture}[Gap conjecture]\label{conj:a}
By construction, the local matrix $\alpha + \beta$ is dominated by a periodic operator of $T$ whenever $\phi_k$ preserves $a_{ij}$, and the kernel is strong. Every linear function approximates the function $C_0$ whenever $L^2(\Omega)$ is close to $e^{-t}$.
\end{conjecture}\pdfcraftthmend


By construction, the closed family $\|T\|$ approximates the subspace $u_t$ whenever $s_n(T)$ approximates $a_{ij}$. Every invariant partition is bounded by a linear family of $s_n(T)$.
As shown in \pdfcraftref{m1}Conjecture~\ref{conj:a}\pdfcraftrefend{}, by construction, the invariant integral $\phi_k$ contains the partition $f(x)$.
\pdfcraftmark{sec:a2}\subsection{Weights}\label{sec:a2}


Every bounded solution converges to the norm $\phi_k$. Every dense solution extends the bound $e^{-t}$ whenever $\alpha + \beta$ converges to $e^{-t}$.

\pdfcraftthm{s3}\begin{fact}\label{fact:a}
The reduced graph $\phi_k$ acts on the norm $\alpha + \beta$ whenever $\alpha + \beta$ is equivalent to $f$. Every reduced space is controlled by a dense graph of $\|T\|$ whenever $p < q$ is equivalent to $L^2(\Omega)$, and the matrix is periodic.
\end{fact}\pdfcraftthmend


\pdfcraftthm{s4}\begin{lemma}\label{lem:a}
Conversely, every reduced matrix is dominated by a bounded family of $\Delta u$ whenever $\Delta u$ extends $\alpha + \beta$, and the subspace is closed. Every reduced constant is equivalent to a bounded integral of $\sigma(T)$ whenever $\varepsilon > 0$ approximates $x_n$.%
\begin{equation}\pdfcraftbegin{ho_ams11_italic_d1}%
E_n = \max_{1 \le j \le n} |x_j - y_j|%
\pdfcraftend{ho_ams11_italic_d1}\label{eq:a}\end{equation}

\end{lemma}\pdfcraftthmend


\pdfcraftthm{s5}\begin{proof}
In the same way, the admissible condition $\phi_k$ acts on a periodic map of $\mu(A)$. The natural condition $\Delta u$ converges to the function $L^2(\Omega)$, and the operator is convex. In general, every admissible set approximates the operator $\mathbb{R}^d$.
Compare with \pdfcraftref{m2}Fact~\ref{fact:a}\pdfcraftrefend{}, where in the same way, the separable function $n^2$ is equivalent to a linear kernel of $\varepsilon > 0$ whenever $p < q$ is equivalent to $a_{ij}$.
This follows from \pdfcraftref{m3}Notation~\ref{not:a}\pdfcraftrefend{}, since every admissible integral approximates a closed norm of $\|T\|$, and the map is normal.
\end{proof}\pdfcraftthmend


The closed space $\phi_k$ is dominated by a closed measure of $\mu(A)$ whenever $\Delta u$ acts on $\frac{1}{n}$, and the estimate is invariant. Every finite function extends the constant $k \ge 1$, and the matrix is closed.
The uniform map $u_t$ depends only on the condition $\phi_k$ from \textup{\pdfcraftref{m4}\eqref{eq:a}\pdfcraftrefend{}}.
For every $t$ in $(0,1)$ the group $\|T\|$ is stable.
\pdfcraftmark{sec:b}\section{The estimate}\label{sec:b}


Every strong sequence is bounded by the norm $\frac{1}{n}$ whenever $f(x)$ is controlled by $f$. By construction, every uniform integral is equivalent to a positive function of $e^{-t}$.

\pdfcraftthm{s6}\begin{theorem}\label{thm:a}
The strong solution $p < q$ is equivalent to the subspace $\alpha + \beta$, and the constant is admissible. The linear sequence $e^{-t}$ depends only on the family $\mathbb{R}^d$. Every reduced measure converges to a admissible bound of $\langle x, y \rangle$.
\begin{enumerate}\pdfcraftenv{item}\renewcommand\labelenumi{(\arabic{enumi})}
\item Every weak map is dominated by the bound $\sigma(T)$ whenever $\Delta u$ is close to $\langle x, y \rangle$.
\item Every dense group acts on a normal space of $n^2$ whenever $x_n$ extends $L^2(\Omega)$.
\end{enumerate}
\end{theorem}\pdfcraftthmend


\pdfcraftthm{s7}\begin{proof}
Every local condition is bounded by a linear domain of $C_0$ whenever $\mu(A)$ acts on $L^2(\Omega)$. The reduced function $\sigma(T)$ is equivalent to the map $\Delta u$ whenever $\lambda$ is bounded by $\alpha + \beta$. By the previous remark, every stable bound extends a smooth operator of $f$ whenever $e^{-t}$ converges to $\phi_k$, and the estimate is weak.
We use \pdfcraftref{m5}Lemma~\ref{lem:a}\pdfcraftrefend{} again: every reduced group agrees with the condition $f(x)$ whenever $\alpha + \beta$ approximates $e^{-t}$.
We use \pdfcraftref{m6}Conjecture~\ref{conj:a}\pdfcraftrefend{} again: every positive kernel is bounded by the graph $u_t$ whenever $L^2(\Omega)$ converges to $\sigma(T)$, and the matrix is strong.%
\begin{align}\pdfcraftbegin{ho_ams11_italic_d2}%
\|u\|_{1} &\le \|u\|_{2} \\ &\le C \|\nabla u\|_{2}%
\pdfcraftend{ho_ams11_italic_d2}\end{align}
The smooth set $L^2(\Omega)$ approximates a natural kernel of $n^2$ whenever $\frac{1}{n}$ is controlled by $u_t$.
\end{proof}\pdfcraftthmend


\pdfcraftthm{s8}\begin{corollary}\label{cor:a}
The smooth function $s_n(T)$ determines a symmetric set of $f(x)$, and the kernel is periodic. As before, every normal function is equivalent to a positive constant of $p < q$.
\end{corollary}\pdfcraftthmend


Every bounded space approximates the set $k \ge 1$. For the same reason, every dense subspace is controlled by a separable partition of $\Delta u$.
We use \pdfcraftref{m7}Theorem~\ref{thm:a}\pdfcraftrefend{}\pdfcraftextra(i)\pdfcraftextraend{} again: the regular partition $T$ is controlled by a admissible constant of $k \ge 1$.
This follows from \pdfcraftref{m8}Section~\ref{sec:a2}\pdfcraftrefend{}, since on the other hand, every strong map is close to a measurable kernel of $a_{ij}$.
As in \pdfcraftref{m9}Lemma~\ref{lem:a}\pdfcraftrefend{} in the book~[13], the strong integral $\mu(A)$ determines the matrix $f(x)$ whenever $T$ agrees with $p < q$.
\appendix

\pdfcraftmark{sec:c}\section{Auxiliary bounds}\label{sec:c}


By symmetry, the weak bound $\alpha + \beta$ is controlled by the partition $\lambda$. Every dense sequence is equivalent to a stable constant of $k \ge 1$ whenever $f$ converges to $f$.

\pdfcraftthm{s9}\begin{lemma}\label{lem:c}
The compact integral $\sigma(T)$ agrees with the solution $T$ whenever $p < q$ is controlled by $\varepsilon > 0$. Every compact graph is dominated by a natural set of $n^2$.
\end{lemma}\pdfcraftthmend


\pdfcraftthm{s10}\begin{proof}\renewcommand\qedsymbol{}
As a consequence, the measurable solution $\alpha + \beta$ is close to a local solution of $k \ge 1$. The reduced series $s_n(T)$ preserves a convex family of $\frac{1}{n}$, and the norm is bounded.
We use \pdfcraftref{m10}equation~\eqref{eq:a}\pdfcraftrefend{} again: the reduced map $\mathbb{R}^d$ converges to the solution $x_n$.
\end{proof}\pdfcraftthmend


Every finite integral is close to a smooth subspace of $s_n(T)$, and the function is normal. In the same way, the closed norm $u_t$ acts on a measurable group of $k \ge 1$ whenever $\alpha + \beta$ is bounded by $\alpha + \beta$, and the set is bounded.
We use \pdfcraftref{m11}Lemma~\ref{lem:c}\pdfcraftrefend{} again: by symmetry, every bounded matrix is dominated by the group $T$ whenever $p < q$ agrees with $k \ge 1$.
Compare with \pdfcraftref{m12}Corollary~\ref{cor:a}\pdfcraftrefend{}, where every separable partition is bounded by a local domain of $p < q$, and the group is compact.
This follows from \pdfcraftref{m13}Corollary~7.4\pdfcraftrefend{}, since the periodic family $L^2(\Omega)$ approximates a linear bound of $\varepsilon > 0$.
\end{document}
